% Caderno 4 (e caderno 2, seção de ortogonalidade) — a grade 2^2 com ponto
% central, no espaço natural e no espaço codificado.
%
% Dados: UFSC, `Cap4_Planejamento_Experimentos_Parte2.pdf`, p. 19 a 21.
% x1 = concentração de catalisador, de 10 a 20; x2 = temperatura, de 20 a 50.
% Ponto central em (15, 35). A soma dos produtos vale 2625 em unidades
% naturais e zero em unidades codificadas: é essa a razão de existir da
% codificação em +-1.
\documentclass[border=3pt,11pt]{standalone}
% Preâmbulo comum das figuras TikZ das apostilas.
% Carregado por todos os `standalone` deste diretório, para que os cinco
% diagramas não divirjam em fonte, cor ou espessura de traço.
%
% fontenc T1 + inputenc utf8 são obrigatórios: sem eles "ç", "ã" e "é"
% saem trocados ou desaparecem silenciosamente do PDF.
\usepackage[T1]{fontenc}
\usepackage[utf8]{inputenc}
\usepackage[brazil]{babel}
% A ordem importa: `amssymb` antes de `newtxmath`, exatamente como em
% `doeapostila.cls`. Invertida, as duas definem `\Bbbk` e a compilação aborta.
\usepackage{amsmath,amssymb}
\usepackage{newtxtext,newtxmath}   % mesma família serifada da classe doeapostila
\usepackage{tikz}
\usetikzlibrary{arrows.meta,calc,decorations.pathreplacing,positioning,
  shapes.geometric,backgrounds}

% Paleta: preto, cinzas e uma única cor de apoio.
\definecolor{azulDoE}{HTML}{1B4E79}
\definecolor{cinzaDoE}{gray}{0.55}
\definecolor{cinzaClaro}{gray}{0.82}
\definecolor{cinzaFundo}{gray}{0.93}

\tikzset{
  % vértice de plano fatorial, com a resposta escrita dentro
  vertice/.style={circle, draw=black, fill=white, line width=0.5pt,
                  inner sep=0pt, minimum size=8.5mm, font=\scriptsize},
  vertdestaque/.style={circle, draw=azulDoE, fill=azulDoE!10, text=azulDoE,
                  line width=0.9pt, inner sep=0pt, minimum size=8.5mm,
                  font=\scriptsize},
  % ponto de planejamento, sem texto dentro
  ponto/.style={circle, draw=black, fill=black, inner sep=0pt,
                minimum size=2.4mm},
  pontovazio/.style={circle, draw=black, fill=white, line width=0.5pt,
                inner sep=0pt, minimum size=2.6mm},
  pontoazul/.style={circle, draw=azulDoE, fill=azulDoE, inner sep=0pt,
                minimum size=2.8mm},
  axial/.style={rectangle, draw=azulDoE, fill=azulDoE, inner sep=0pt,
                minimum size=2.6mm},
  apagado/.style={circle, draw=cinzaClaro, fill=white, line width=0.5pt,
                inner sep=0pt, minimum size=2.6mm},
  aresta/.style={draw=black, line width=0.5pt},
  arestafundo/.style={draw=cinzaClaro, line width=0.5pt},
  eixo/.style={draw=black, line width=0.5pt, -{Latex[length=1.8mm,width=1.4mm]}},
  seta/.style={draw=black, line width=0.6pt, -{Latex[length=2.2mm,width=1.7mm]}},
  setaazul/.style={draw=azulDoE, line width=0.8pt,
                   -{Latex[length=2.2mm,width=1.7mm]}},
  cota/.style={draw=cinzaDoE, line width=0.5pt,
               {Latex[length=1.6mm]}-{Latex[length=1.6mm]}},
  % Corpos escolhidos para casar com as figuras de matplotlib, que usam
  % 8 pt nos ticks e 9 pt nos rótulos de eixo. Com a base de 11 pt do
  % `standalone` (a mesma da classe doeapostila), \scriptsize = 8 pt e
  % \footnotesize = 9 pt. Abaixo disso a figura deixa de ser legível na
  % página impressa, e é para isso que a folha de prova serve.
  rotulo/.style={font=\footnotesize},
  rotulinho/.style={font=\scriptsize},
  nota/.style={font=\footnotesize, align=left},
  notaazul/.style={font=\footnotesize, align=left, text=azulDoE},
}

\begin{document}
\begin{tikzpicture}[x=1cm, y=1cm]

  % ======================= painel (a): natural ==========================
  \begin{scope}
    \draw[eixo] (0,0) -- (4.5,0) node[rotulo, right, yshift=-0.5mm] {$x_1$};
    \draw[eixo] (0,0) -- (0,4.5) node[rotulo, above] {$x_2$};

    \coordinate (a11) at (1.1,1.1);
    \coordinate (a21) at (3.7,1.1);
    \coordinate (a12) at (1.1,3.7);
    \coordinate (a22) at (3.7,3.7);
    \coordinate (a00) at (2.4,2.4);

    \draw[cinzaClaro, line width=0.5pt] (a11) -- (a21) -- (a22) -- (a12) -- cycle;
    \draw[cinzaClaro, dotted] (a00) -- (1.1,2.4);
    \draw[cinzaClaro, dotted] (a00) -- (2.4,1.1);

    \node[pontovazio] at (a11) {};
    \node[pontovazio] at (a21) {};
    \node[pontovazio] at (a12) {};
    \node[pontovazio] at (a22) {};
    \node[pontoazul]  at (a00) {};
    \node[rotulinho, text=azulDoE, anchor=south west] at (2.52,2.46)
      {ponto central};

    \draw[cinzaDoE, line width=0.4pt] (1.1,0.06) -- (1.1,-0.06);
    \draw[cinzaDoE, line width=0.4pt] (2.4,0.06) -- (2.4,-0.06);
    \draw[cinzaDoE, line width=0.4pt] (3.7,0.06) -- (3.7,-0.06);
    \node[rotulinho, below] at (1.1,-0.1) {10};
    \node[rotulinho, below] at (2.4,-0.1) {15};
    \node[rotulinho, below] at (3.7,-0.1) {20};

    \draw[cinzaDoE, line width=0.4pt] (0.06,1.1) -- (-0.06,1.1);
    \draw[cinzaDoE, line width=0.4pt] (0.06,2.4) -- (-0.06,2.4);
    \draw[cinzaDoE, line width=0.4pt] (0.06,3.7) -- (-0.06,3.7);
    \node[rotulinho, left] at (-0.1,1.1) {20};
    \node[rotulinho, left] at (-0.1,2.4) {35};
    \node[rotulinho, left] at (-0.1,3.7) {50};

    \node[nota, anchor=north] at (2.25,-0.75)
      {(a) unidades naturais};
    \node[rotulinho, anchor=north, align=center] at (2.25,-1.25)
      {$x_1$: concentração de catalisador\\ $x_2$: temperatura};
    \node[nota, anchor=north] at (2.25,-2.35)
      {$\displaystyle\sum_k x_{1k}x_{2k} = 2625 \neq 0$};
  \end{scope}

  % ============== seta central com a fórmula de codificação =============
  \begin{scope}[xshift=5.2cm]
    \draw[seta] (0,2.3) -- (2.5,2.3);
    \node[nota, anchor=south, align=center] at (1.25,2.5)
      {$z_i = \dfrac{x_i - \dfrac{x_i^{\min}+x_i^{\max}}{2}}
                    {\dfrac{x_i^{\max}-x_i^{\min}}{2}}$};
    \node[rotulinho, anchor=north, align=center] at (1.25,2.1)
      {codificação};
  \end{scope}

  % ======================= painel (b): codificado =======================
  \begin{scope}[xshift=9.4cm]
    \draw[eixo] (0,2.4) -- (4.5,2.4) node[rotulo, right, yshift=-0.5mm] {$z_1$};
    \draw[eixo] (2.4,0) -- (2.4,4.5) node[rotulo, above] {$z_2$};

    \coordinate (b11) at (1.1,1.1);
    \coordinate (b21) at (3.7,1.1);
    \coordinate (b12) at (1.1,3.7);
    \coordinate (b22) at (3.7,3.7);
    \coordinate (b00) at (2.4,2.4);

    \draw[cinzaClaro, line width=0.5pt] (b11) -- (b21) -- (b22) -- (b12) -- cycle;

    \node[pontovazio] at (b11) {};
    \node[pontovazio] at (b21) {};
    \node[pontovazio] at (b12) {};
    \node[pontovazio] at (b22) {};
    \node[pontoazul]  at (b00) {};

    \node[rotulinho, below left]  at (b11) {$(-1,-1)$};
    \node[rotulinho, below right] at (b21) {$(+1,-1)$};
    \node[rotulinho, above left]  at (b12) {$(-1,+1)$};
    \node[rotulinho, above right] at (b22) {$(+1,+1)$};
    \node[rotulinho, text=azulDoE, anchor=south west] at (2.52,2.46)
      {$(0,0)$};

    \draw[cinzaDoE, line width=0.4pt] (1.1,2.46) -- (1.1,2.34);
    \draw[cinzaDoE, line width=0.4pt] (3.7,2.46) -- (3.7,2.34);
    \node[rotulinho, below] at (1.1,2.30) {$-1$};
    \node[rotulinho, below] at (3.7,2.30) {$+1$};
    \draw[cinzaDoE, line width=0.4pt] (2.46,1.1) -- (2.34,1.1);
    \draw[cinzaDoE, line width=0.4pt] (2.46,3.7) -- (2.34,3.7);
    \node[rotulinho, left] at (2.30,1.1) {$-1$};
    \node[rotulinho, left] at (2.30,3.7) {$+1$};

    \node[nota, anchor=north] at (2.25,-0.75)
      {(b) unidades codificadas};
    \node[rotulinho, anchor=north, align=center] at (2.25,-1.25)
      {a origem passa a ser o ponto central\\ e a grade fica simétrica};
    \node[notaazul, anchor=north] at (2.25,-2.35)
      {$\displaystyle\sum_k z_{1k}z_{2k} = 0$};
  \end{scope}

\end{tikzpicture}
\end{document}
